\documentclass[12pt]{article}
\usepackage{graphicx} % Required for inserting images
\usepackage[margin=0.5in]{geometry}
\usepackage{amsmath, amssymb}
\usepackage{mathrsfs}


\usepackage{soul}


% Dr. David Brown Commands
\newcommand{\ds}{\displaystyle}
\newcommand{\R}{\mathbb{R}}
\newcommand{\N}{\mathbb{N}}
\newcommand{\Q}{\mathbb{Q}}
\newcommand{\C}{\mathbb{C}}


% Commands to get script r
\usepackage{calligra}
\DeclareMathAlphabet{\mathcalligra}{T1}{calligra}{m}{n}
\DeclareFontShape{T1}{calligra}{m}{n}{<->s*[2.2]callig15}{}
\newcommand{\scripty}[1]{\ensuremath{\mathcalligra{#1}}}

% Unit commands
\newcommand{\degree}{\text{°}}
\newcommand{\unitSpeed}{\frac{\text{m}}{\text{s}}}
\newcommand{\unitAcceleration}{\frac{\text{m}}{\text{s}^2}}

% Constant commands
\newcommand{\avogadro}{6.022\times10^{23}}

\newcommand{\boltzmannConstK}{1.381\times10^{-23}\frac{\text{J}}{\text{K}}}
\newcommand{\boltzmannConstInEV}{8.617\times10^{-5}\frac{\text{eV}}{\text{K}}}
\newcommand{\planckConst}{6.626\times10^{-34}\text{J}\cdot\text{s}}

\newcommand{\atmP}{1.01325\times10^5\,\text{Pa}}

% Known Value Commands
\newcommand{\electronMass}{9.109\times10^{-31}\,\text{kg}}
\newcommand{\protonMass}{1.673\times10^{-27}\,\text{kg}}

% Commands to easily write partial derivatives
\newcommand{\partDeriv}[2]{\frac{\partial #1}{\partial #2}}
\newcommand{\doublePartDeriv}[2]{\frac{\partial^2 #1}{\partial^2 #2}}


\title{Problem 2.4}
\author{Gabriel Decker}


\begin{document}


\maketitle
\begin{flushleft}


In this problem, we find $\vec{E}$ a distance $z$ above the center of a square loop of charge with uniform density $\lambda$ and side length $a$, according to Figure 2.8.
Let's define the coordinate system with the origin at the center of the charge square and square charge being in the $xy$-plane and the point of interest extending in the $\hat{z}$ direction.
So, $\hat{r}=z\hat{z}$.
It's easy to see by the symmetry of the problem that the $\hat{x}$ and $\hat{y}$ components of the electric field at the point of interest is zero because each side has identical charge, so a bit of charge on one side whould have the $\hat{x}$ and $\hat{y}$ components of its electric field cancelled out by a bit of charge on the other side of the square.
So, we only need to find the $\hat{z}$ component of the electric field from each side.
Since every side is the same and is the same distance from the center, this implies that the $\hat{z}$ component of the electric field, $E_z$ is four times the $\hat{z}$ component from each side $E_{za}$, so $E_z=4E_{za}$.
So, if we can get just the $\hat{z}$ component of the electric field from one side of the charge square, then this solves our problem.

As directed, we can use the following equation derived in example 2.2.
The electric field a distance $z$ above the midpoint of a straight line segment of length $2L$ is the following.
$$\vec{E}=\frac{1}{4\pi\epsilon_0}\frac{2\lambda L}{z\sqrt{z^2+L^2}}\hat{z}$$
where $\hat{z}$ refers to the direction pointing from the center of the line segment to the point of interest.
Recall that in our coordinate system, this could include some $\hat{x}$, $\hat{y}$, and $\hat{z}$ components.
So, we need just the $\hat{z}$ component in our current coordinate system.
Also, recall that in our problem, the side length is $a$, so $L=a/2$, making our version of the formula the following (still in the example's coordinate system).
$$\vec{E}=\frac{1}{4\pi\epsilon_0}\frac{\lambda a}{z\sqrt{z^2+(\frac{a}{2})^2}}\hat{z}$$
Also, the $z$ in the example is the shortest distance between the center of the segment and the point of interest, while in our problem, $z$ goes from the center of the square to the point.
So, $z$ becomes $\sqrt{(\frac{a}{2})^2+z^2}$ by the pythagorean theorem, making our version of the formula the following (still in the example's coordinate system).
$$\vec{E}=\frac{1}{4\pi\epsilon_0}\frac{\lambda a}{\sqrt{(\frac{a}{2})^2+z^2}\sqrt{(\frac{a}{2})^2+z^2+(\frac{a}{2})^2}}\hat{z}$$\\~\\

If we define the angle $\theta$ as the angle between the $xy$-plane and the vector pointing from the center of the line segment to the point of interest, then this implies that $E_{za}=E\sin(\theta)$.
$\ds\sin(\theta)=\frac{z}{\sqrt{(\frac{a}{2})^2+z^2}}$, so $\ds E_{za}=\frac{z}{\sqrt{(\frac{a}{2})^2+z^2}}E$, and thus $\ds E_z=\frac{4z}{\sqrt{(\frac{a}{2})^2+z^2}}E$.\\~\\

Since $E$ is the magnitude of $\vec{E}$, which only has one component, this implies that the $\hat{z}$ component of $\vec{E}$ is the following.
$$E_z=\frac{1}{4\pi\epsilon_0}\frac{4\lambda az}{\left((\frac{a}{2})^2+z^2\right)\sqrt{2(\frac{a}{2})^2+z^2}}$$\\~\\

And thus, the electric field above the center of a square line charge is the following.
$$\vec{E}=\frac{1}{4\pi\epsilon_0}\frac{4\lambda az}{\left((\frac{a}{2})^2+z^2\right)\sqrt{2(\frac{a}{2})^2+z^2}}\hat{z}$$



\end{flushleft}

\end{document}
